% =========================================================================
% DRAFT NOTES: Research Questions, Critique, and Refined Proposals
% This block is commented out using % so it won't affect PDF compilation.
% You can read it here, discuss with your advisor, and write your introduction.
% =========================================================================
%
% -------------------------------------------------------------------------
% 1. ALKUPERÄINEN SUUNNITELMA (Approved Research Plan RQs - 3 menetelmää)
% -------------------------------------------------------------------------
% RQ1: To what extent (in percentage) do small-to-big retrieval, semantic 
%      re-ranking, and prompt caching decrease the Time to First Token (TTFT) 
%      to improve latency?
% RQ2: To what extent (in percentage) do small-to-big retrieval, semantic 
%      re-ranking, and prompt caching decrease the number of tokens spent 
%      per LLM call to reduce costs?
% RQ3: To what extent do small-to-big retrieval, semantic re-ranking, and 
%      prompt caching affect the quality of the output (measured via, for 
%      example, Faithfulness, Answer Relevance, and/or Context Precision) 
%      versus the base case?
% RQ4: (Optional) Are there observable differences in efficiency between 
%      the platforms of different cloud providers?
%
% -------------------------------------------------------------------------
% 2. KRITIIKKI JA HUOMIOT (Critique of the Current RQs)
% -------------------------------------------------------------------------
% * RQ1 (Latenssi/TTFT): 
%   - TTFT (Time to First Token) on loistava interaktiivisuuden mittari, mutta 
%     se ei kerro koko totuutta.
%   - Uudelleenjärjestely (Semantic Re-ranking) ja reititys (Model Cascading) 
%     vaativat erillisiä laskentavaiheita ennen vastaamisen aloitusta (overhead). 
%     Tämä voi *kasvattaa* TTFT:tä naiveen RAG:iin verrattuna, vaikka kokonaistehokkuus 
%     parantuisi.
%   - Pitäisi mitata latenssiprofiilia laajemmin: sekä TTFT:tä että kokonaislatenssia 
%     (End-to-End Latency).
%
% * RQ2 (Kustannukset/Tokenit):
%   - Tokeneiden määrän mittaaminen prosentteina on harhaanjohtavaa, koska:
%     1. Prompt Caching (kontekstin välimuisti mallitasolla) ei vähennä token-määrää, 
%        vaan antaa alennuksen (esim. -50% - -75%) tallennetuille tokeneille. 
%        Pelkkä token-määrän arviointi piilottaa tämän säästön.
%     2. Model Cascading (reititys halvemmalle mallille, esim. Gemini Flash) 
%        ei välttämättä vähennä tokeneita, mutta hinta per token on moninkertaisesti 
%        halvempi (esim. Gemini Flash vs Pro hintaero).
%   - Mittarin tulee mitata todellisia operatiivisia kustannuksia (Operational Cost / $) 
%     pelkän token-määrän sijaan.
%
% * RQ3 (Laatu/Triadi):
%   - RAG-triadi (Faithfulness, Answer Relevance, Context Precision) on standardi.
%   - Reititys (Cascading) tuo uuden riskin: jos reititin ohjaa kyselyn No-RAG -polulle 
%     tai liian pienelle mallille virheellisesti, laatu heikkenee huomattavasti. 
%     Miten optimointi vaikuttaa eri laatu-ulottuvuuksiin erikseen (esim. re-ranking 
%     parantaa Context Precisionia, kun taas small-to-big parantaa Context Recallia).
%
% * RQ4 (Pilvialustat vs. Yhteisvaikutus):
%   - Pilvialustojen vertailu (GCP vs. muut) siirtää fokuksen pois RAG-algoritmeista 
%     ja vaatii suuren infrastruktuurityön.
%   - Enterprise-ympäristössä kriittisempää on tietää optimointimenetelmien 
%     *yhteisvaikutus* (Stacking / Combinations). Esim. nollaako Semantic Re-ranking 
%     Prompt Cachingin hyödyt muuttamalla promptia liikaa? Mikä on optimaalinen 
%     tasapaino (Pareto-optimi) kustannusten, latenssin ja laadun välillä?
%
% -------------------------------------------------------------------------
% 3. EHDOTUKSET UUSIKSI TUTKIMUSKYSYMYKSIKSI (Refined Research Questions)
% -------------------------------------------------------------------------
%
% Taso A: Yksinkertainen laajennus neljälle menetelmälle (jos halutaan säilyttää
% alkuperäisen suunnitelman rakenne ja henki):
%
% - RQ1: To what extent (in percentage) do small-to-big retrieval, semantic 
%        re-ranking, prompt caching, and model cascading decrease the Time to 
%        First Token (TTFT) to improve latency?
% - RQ2: To what extent (in percentage) do small-to-big retrieval, semantic 
%        re-ranking, prompt caching, and model cascading decrease the number of 
%        tokens spent per LLM call to reduce costs?
% - RQ3: To what extent do small-to-big retrieval, semantic re-ranking, prompt 
%        caching, and model cascading affect the quality of the output 
%        (measured via, for example, Faithfulness, Answer Relevance, and/or 
%        Context Precision) versus the base case?
% - RQ4: (Optional) Are there observable differences in efficiency between the 
%        platforms of different cloud providers?
%
% Taso B (Suositeltu): Kriittisesti parannellut tutkimuskysymykset varsinaiseen työhön:
%
% - RQ1: How do small-to-big retrieval, semantic re-ranking, prompt caching, and 
%        model cascading affect the latency profile (specifically Time to First 
%        Token [TTFT] and overall end-to-end execution latency) of a RAG system?
% - RQ2: To what extent do these optimization methods reduce the operational cost 
%        of the system, accounting for token pricing tiers, prompt caching 
%        discounts, and model cascading?
% - RQ3: How do these optimizations impact retrieval and generation quality 
%        (measured by Context Precision, Context Recall, Faithfulness, and 
%        Answer Relevance) compared to a naive RAG baseline?
% - RQ4: How do these optimization methods interact when combined, and what is 
%        the optimal equilibrium (Pareto frontier) between latency, cost, and 
%        quality in an enterprise customer support system?
% =========================================================================

Modern enterprise-grade AI systems face continuously increasing resource constraints, in addition to growing operational expenditures, driven by a rapid acceleration in corporate AI integration and exponential growth in the demand for AI computing capacity \citep{StanfordAIIndex2026}. Although API providers frequently lower unit token pricing, overall costs continue to rise due to explosive token usage and the complexity of multi-step agentic workflows that uses the latest foundational large-language models (LLMs) \citep{chen2023frugalgpt, GoogleVertexAIPricing}. This highlights the critical need for optimization across all phases of enterprise AI systems to ensure that computing resources are utilized as efficiently as possible.

This is especially true for retrieval-augmented generation (RAG), which is one of the most prominent solutions to one of the most fundamental drawbacks of the modern LLMs - hallucination. Due to its architecture, it increases the token, especially input, consumption as the core idea of the RAG is to provide additional context for the LLM to provide more accurate answers. However, many studies have presented optimization strategies that are tested on pre-defined data sets and are not implemented using tools that are available for industry standard tool, such as in SpringAI. 

The aim of the thesis is to research what are the state-of-the-art optimization strategies and the main contribution is to test how does these strategies work in enterprise environments. In addition, it will provide one possible solution to implement these strategies using the enterprise frameworks. To limit the scope of the thesis, four different strategies are taken into deeper consideration. The methods chosen for the thesis are cascading, small-to-big retrieval, semantic re-ranking, and prompt caching. Each of these methods takes a different viewpoint on the RAG process. Cascading and small-to-big retrieval is a method for optimizing the query phase of the pipeline, whereas semantic re-ranking is used to re-rank the query results to optimize the final prompt passed to the LLM. Prompt caching, on the other hand, is a method to store results in a cache so that similar questions can receive answers directly, without driving the whole computationally heavy RAG pipeline.

As the enterprise systems must manage costs, have low latency and answers have good quality, these factors are taken into consideration. 

<table here>

%% 1. Describe the overall context + bring several background references
%% 2. Show the research gap (why to do this particular research)
%% 3. Bring up your research question(s)
%% 4. Describe the structure of the work + mention research method
%% 
%% Reference example: \cite{travers1996programming}
%% 
%% Research question: Fill this up
%% Research methods: survey / literature review
